\normalsize

\chapter{引言}

\section{研究背景}

随着边缘智能、物联网和计算机视觉技术的发展，视频监控系统正在由传统的被动录像设备逐步演进为具备实时感知、边缘分析和主动控制能力的智能终端。在早期视频监控应用中，摄像机的主要作用是采集和存储画面，后续分析通常依赖人工查看或中心平台回放。随着场景复杂度和响应实时性要求的提高，监控设备不仅需要“看得见”，还需要能够在本地完成目标识别、事件判断和运动控制，从而缩短从感知到执行的链路。

布控球作为一种便携式移动监控设备，集成摄像头、云台、通信模块和供电模块，能够快速部署在临时布防、应急救援、施工巡检、交通保障和偏远区域监测等场景中。相较固定摄像机，布控球具备部署灵活、覆盖范围可调和远程操控能力强等优势，因此在智慧城市和行业安防系统中具有较高的工程价值。在实际使用中，布控球往往被安装在三脚架、车辆、杆体或临时支架上，设备姿态和监控视角会随部署条件变化而变化，这使得其自动化控制能力直接影响现场使用效率。

然而，传统布控球的云台控制方式仍以人工遥控、预置点切换或固定轨迹巡航为主。当画面中出现快速移动目标、多人同时进入、目标短时遮挡或光照突变时，人工操作难以及时保持目标居中，固定巡航策略也无法根据目标位置动态调整云台姿态。若将视频流全部上传云端处理，系统又会受到网络带宽、传输时延和隐私安全等因素限制。在应急处置和移动巡检场景下，网络条件并不总是稳定，现场人员也难以长期保持高强度监控，因此仅依赖人工和云端的方案难以满足低延迟、低成本和高可用性的综合要求。

因此，将轻量级视觉检测算法和自适应云台控制策略部署在边缘端，使布控球能够在本地完成目标检测、目标选择和控制决策，是提升移动监控设备智能化水平的重要方向。边缘端方案的核心价值在于将“发现目标”和“调整视角”两个环节压缩在设备附近完成：一方面，系统能够减少视频全量上传带来的带宽压力；另一方面，控制指令可以根据最新画面快速生成，有利于提高目标居中和持续跟踪的稳定性。

本课题面向海康威视布控球设备，围绕“视频采集--目标检测--目标记忆--云台控制--人机交互”的闭环链路，研究并实现一种自适应云台控制嵌入式设备原型。系统以 Python 为主要开发语言，集成 HCNetSDK、PlayCtrl、YOLOv8n、OpenCV、Tkinter 和轻量级人体 ReID 模块，实现实时视频流获取、人体目标检测、目标编号与选择、PTZ 自动跟踪、倒挂模式适配以及手动控制等功能。与单纯的离线检测算法验证不同，本课题强调真实设备接入和系统闭环运行，重点解决 SDK 调用、视频流解码、检测线程调度、云台控制节奏和界面交互之间的协同问题。

\section{研究意义}

本研究具有一定的工程应用意义。首先，系统将目标检测与 PTZ 控制部署在边缘端，减少了对云端实时分析链路的依赖，有利于降低网络传输压力并提升控制响应速度。其次，系统面向实际布控球设备完成 SDK 调用、码流解码、实时帧获取和云台控制接口适配，能够为同类设备的二次开发提供可复用经验。再次，本文在简单最大目标跟踪的基础上加入人体 ReID 记忆和指定 ID 跟踪机制，使系统在多人场景下具备更好的目标连续性，避免频繁切换跟踪对象。最后，系统提供可视化控制界面和倒挂模式切换功能，使设备更贴近实际部署和调试需求。

从项目实现过程看，本研究还具有一定的综合训练价值。课题同时涉及设备通信、接口适配、实时图像处理、深度学习推理、目标关联、反馈控制和人机交互设计等内容。各模块单独实现并不复杂，但当它们被放入同一个实时系统中时，就需要考虑数据格式、线程节奏、异常恢复和资源释放等工程问题。因此，该课题能够较好地体现人工智能算法从模型调用走向设备应用时所必须面对的系统集成过程。

从学术和技术融合角度看，本课题将轻量级目标检测、视觉伺服控制、嵌入式实时处理和视频监控设备接口适配结合起来。它既关注算法在边缘环境下的实时性，也关注云台控制的稳定性和人机交互的可用性，对计算机视觉算法工程化落地具有一定参考价值。本文并不追求提出新的通用检测网络，而是围绕布控球这一具体设备，探索如何将成熟模型、轻量目标记忆和规则控制策略组合成可运行的应用系统。

\section{研究目标}

结合任务书、开题报告和中期检查材料，本文的研究目标包括以下几个方面：

\begin{enumerate}
    \item 完成布控球硬件接口适配。基于海康威视 HCNetSDK 与 PlayCtrl 播放库，实现设备初始化、登录认证、实时预览、码流解码和 PTZ 控制等功能，构建稳定的视频采集与控制通道。
    \item 实现轻量级人体目标检测模块。选用 YOLOv8n 作为检测模型，完成实时帧读取、人体类别筛选、置信度过滤、目标坐标输出和检测结果可视化，为后续控制算法提供输入。
    \item 设计自适应云台控制策略。根据目标中心与画面中心偏差生成 PTZ 控制指令，并结合阈值死区、比例脉冲、指令冷却和指数平滑机制降低跟踪抖动。
    \item 增强多目标场景下的目标连续性。构建轻量级人体 ReID 记忆模块，为人体目标分配稳定 ID，支持自动目标选择和指定 ID 跟踪。
    \item 完成系统集成与测试。将视频采集、目标检测、目标记忆、云台控制和图形界面整合为完整原型，并从实时性、稳定性和可用性等方面进行验证。
\end{enumerate}

上述目标之间存在明显的递进关系。硬件接口适配是系统运行的基础，只有稳定获得视频帧并能够下发云台指令，后续检测和控制才具备实际意义；目标检测模块提供目标位置，是闭环控制的数据源；人体 ReID 记忆模块用于改善多目标场景下的目标一致性；PTZ 控制策略则把图像坐标偏差转换为可执行的机械动作；最终，图形界面和测试验证用于保证系统能够被用户观察、调试和复现。

\section{主要研究内容}

围绕上述目标，本文开展的主要研究内容如下：

\subsection{硬件接口适配与视频采集}

系统通过 HCNetSDK 完成布控球设备的初始化、登录和云台控制，通过 PlayCtrl 完成实时码流解码，并将解码后的图像帧缓存为检测模块可读取的数据。该部分重点解决 SDK 接口适配、设备登录失败处理、预览句柄管理、解码回调保存和资源释放等工程问题。由于 SDK 回调数据并不是检测模型可以直接处理的图像格式，系统需要在解码回调、图像保存、帧读取和尺寸归一化之间建立稳定的数据通路，并保证在检测线程读取图像时不会阻塞预览链路。

\subsection{轻量级目标检测与坐标输出}

检测模块选用 YOLOv8n 预训练模型，对布控球画面中的人体目标进行实时检测。系统设置输入尺寸、置信度阈值和半精度推理选项，在实时性与检测质量之间取得平衡。检测结果经过类别筛选和坐标转换后，输出边界框、中心点、置信度和面积等信息，为目标选择和 PTZ 控制提供结构化数据。由于布控球画面可能存在视角变化、目标远近变化和背景干扰，系统没有直接使用全部检测类别，而是聚焦人体类别，并将检测框面积作为自动目标选择的重要依据。

\subsection{人体 ReID 记忆与目标选择}

针对多目标场景中“最大目标策略”容易发生目标切换的问题，本文设计了轻量级人体 ReID 记忆模块。该模块利用人体框位置比例和平均颜色等低成本特征计算相似度，在短时缺失后仍可将重新出现的目标合并为同一 ID。系统支持自动跟踪最大人体，也支持用户在界面中循环选择指定 ID。这样既保留了自动模式下快速选择主目标的便利性，又允许用户在多人场景中明确指定需要跟踪的对象，从而提高系统的可控性。

\subsection{自适应 PTZ 控制}

控制模块根据目标中心点与画面中心点的偏差确定水平和垂直转动方向。为避免微小检测抖动引发云台频繁动作，系统设置水平和垂直死区阈值；为兼顾远距离偏差的快速回正与近中心区域的细致调整，系统根据偏差大小调整单次 PTZ 脉冲时长，并引入冷却间隔限制控制指令频率。该策略虽然没有建立严格的云台动力学模型，但参数直观、调试成本低，适合在原型系统和真实设备测试中逐步修正。

\subsection{可视化界面与系统集成}

系统构建人机交互模块，提供视频预览、目标状态呈现、跟踪模式切换、倒挂适配和人工接管等能力。交互层与后台检测线程、PTZ 控制线程协同工作，使测试人员能够观察目标框、目标 ID、偏移量和云台运动效果。该模块的设计重点不是复杂交互形式，而是支撑系统状态管理、控制方向验证和异常场景下的人工干预。

\section{论文结构}

本文共分为六章。第一章介绍课题背景、研究意义、研究目标和主要内容。第二章综述布控球云台控制、目标检测、目标跟踪与边缘智能相关技术。第三章分析系统需求并给出总体架构设计。第四章详细说明视频采集、目标检测、人体 ReID、PTZ 控制和界面集成等关键模块的实现。第五章给出系统测试方案、初步测试结果和问题分析。第六章总结全文工作，并对模型优化、复杂场景适应性和云边协同方向进行展望。
